\documentclass[10pt,a4paper]{amsart}
\usepackage{amsmath,amssymb,amsthm,graphicx}
\usepackage{verbatim}
\newenvironment{reference}{\comment}{\endcomment}
%\newenvironment{reference}{}{}
\newenvironment{slogan}{\comment}{\endcomment}
\newenvironment{history}{\comment}{\endcomment}

% For commutative diagrams we use Xy-pic
\usepackage[all]{xy}

% We use 2cell for 2-commutative diagrams.
\xyoption{2cell}
% Theorem environments.
%
\theoremstyle{plain}
\newtheorem{theorem}[subsection]{Theorem}
\newtheorem{proposition}[subsection]{Proposition}
\newtheorem{lemma}[subsection]{Lemma}

\theoremstyle{definition}
\newtheorem{definition}[subsection]{Definition}
\newtheorem{example}[subsection]{Example}
\newtheorem{exercise}[subsection]{Exercise}
\newtheorem{situation}[subsection]{Situation}

\theoremstyle{remark}
\newtheorem{remark}[subsection]{Remark}
\newtheorem{remarks}[subsection]{Remarks}

\numberwithin{equation}{subsection}

% Macros
%
\def\lim{\mathop{\mathrm{lim}}\nolimits}
\def\colim{\mathop{\mathrm{colim}}\nolimits}
\def\Spec{\mathop{\mathrm{Spec}}}
\def\Hom{\mathop{\mathrm{Hom}}\nolimits}
\def\Ext{\mathop{\mathrm{Ext}}\nolimits}
\def\SheafHom{\mathop{\mathcal{H}\!\mathit{om}}\nolimits}
\def\SheafExt{\mathop{\mathcal{E}\!\mathit{xt}}\nolimits}
\def\Sch{\mathit{Sch}}
\def\Mor{\mathop{\mathrm{Mor}}\nolimits}
\def\Ob{\mathop{\mathrm{Ob}}\nolimits}
\def\Sh{\mathop{\mathit{Sh}}\nolimits}
\def\NL{\mathop{N\!L}\nolimits}
\def\CH{\mathop{\mathrm{CH}}\nolimits}
\def\proetale{{pro\text{-}\acute{e}tale}}
\def\etale{{\acute{e}tale}}
\def\QCoh{\mathit{QCoh}}
\def\Ker{\mathop{\mathrm{Ker}}}
\def\Im{\mathop{\mathrm{Im}}}
\def\Coker{\mathop{\mathrm{Coker}}}
\def\Coim{\mathop{\mathrm{Coim}}}

% Boxtimes
%
\DeclareMathSymbol{\boxtimes}{\mathbin}{AMSa}{"02}

%
% Macros for moduli stacks/spaces
%
\def\QCohstack{\mathcal{QC}\!\mathit{oh}}
\def\Cohstack{\mathcal{C}\!\mathit{oh}}
\def\Spacesstack{\mathcal{S}\!\mathit{paces}}
\def\Quotfunctor{\mathrm{Quot}}
\def\Hilbfunctor{\mathrm{Hilb}}
\def\Curvesstack{\mathcal{C}\!\mathit{urves}}
\def\Polarizedstack{\mathcal{P}\!\mathit{olarized}}
\def\Complexesstack{\mathcal{C}\!\mathit{omplexes}}
% \Pic is the operator that assigns to X its picard group, usage \Pic(X)
% \Picardstack_{X/B} denotes the Picard stack of X over B
% \Picardfunctor_{X/B} denotes the Picard functor of X over B
\def\Pic{\mathop{\mathrm{Pic}}\nolimits}
\def\Picardstack{\mathcal{P}\!\mathit{ic}}
\def\Picardfunctor{\mathrm{Pic}}
\def\Deformationcategory{\mathcal{D}\!\mathit{ef}}
\usepackage{cleveref}
\crefname{originalreference}{original reference}{original references}
\Crefname{originalreference}{Original reference}{Original references}
\makeatletter
\newlabel{situation-setup-base-change}{{0FVG}{source}{}{sourcepointer-1}{}}
\newlabel{situation-setup-base-change@cref}{{[originalreference][0][]0FVG}{[1][0][]source}{}{}{}}
\makeatother
\crefname{originalreference}{original reference}{original references}
\Crefname{originalreference}{Original reference}{Original references}
\makeatletter
\newlabel{situation-setup}{{02QL}{source}{}{sourcepointer-2}{}}
\newlabel{situation-setup@cref}{{[originalreference][0][]02QL}{[1][0][]source}{}{}{}}
\makeatother
\crefname{originallemma}{original lemma}{original lemmas}
\Crefname{originallemma}{Original lemma}{Original lemmas}
\makeatletter
\newlabel{lemma-compose-base-change}{{0FVP}{source}{}{sourcepointer-3}{}}
\newlabel{lemma-compose-base-change@cref}{{[originallemma][0][]0FVP}{[1][0][]source}{}{}{}}
\makeatother
\crefname{originallemma}{original lemma}{original lemmas}
\Crefname{originallemma}{Original lemma}{Original lemmas}
\makeatletter
\newlabel{limits-lemma-descend-finite-presentation}{{01ZM}{source}{}{sourcepointer-4}{}}
\newlabel{limits-lemma-descend-finite-presentation@cref}{{[originallemma][0][]01ZM}{[1][0][]source}{}{}{}}
\makeatother
\crefname{originallemma}{original lemma}{original lemmas}
\Crefname{originallemma}{Original lemma}{Original lemmas}
\makeatletter
\newlabel{limits-lemma-descend-closed-immersion-finite-presentation}{{01ZP}{source}{}{sourcepointer-5}{}}
\newlabel{limits-lemma-descend-closed-immersion-finite-presentation@cref}{{[originallemma][0][]01ZP}{[1][0][]source}{}{}{}}
\makeatother
\crefname{originallemma}{original lemma}{original lemmas}
\Crefname{originallemma}{Original lemma}{Original lemmas}
\makeatletter
\newlabel{lemma-dimension-base-change}{{0FVH}{source}{}{sourcepointer-6}{}}
\newlabel{lemma-dimension-base-change@cref}{{[originallemma][0][]0FVH}{[1][0][]source}{}{}{}}
\makeatother
\crefname{originallemma}{original lemma}{original lemmas}
\Crefname{originallemma}{Original lemma}{Original lemmas}
\makeatletter
\newlabel{limits-lemma-limit-nonempty}{{01Z2}{source}{}{sourcepointer-7}{}}
\newlabel{limits-lemma-limit-nonempty@cref}{{[originallemma][0][]01Z2}{[1][0][]source}{}{}{}}
\makeatother
\crefname{originallemma}{original lemma}{original lemmas}
\Crefname{originallemma}{Original lemma}{Original lemmas}
\makeatletter
\newlabel{limits-lemma-descend-opens}{{01Z4}{source}{}{sourcepointer-8}{}}
\newlabel{limits-lemma-descend-opens@cref}{{[originallemma][0][]01Z4}{[1][0][]source}{}{}{}}
\makeatother
\crefname{originallemma}{original lemma}{original lemmas}
\Crefname{originallemma}{Original lemma}{Original lemmas}
\makeatletter
\newlabel{limits-lemma-descend-section}{{01Z0}{source}{}{sourcepointer-9}{}}
\newlabel{limits-lemma-descend-section@cref}{{[originallemma][0][]01Z0}{[1][0][]source}{}{}{}}
\makeatother
\crefname{originallemma}{original lemma}{original lemmas}
\Crefname{originallemma}{Original lemma}{Original lemmas}
\makeatletter
\newlabel{algebra-lemma-pullback-module}{{02M1}{source}{}{sourcepointer-10}{}}
\newlabel{algebra-lemma-pullback-module@cref}{{[originallemma][0][]02M1}{[1][0][]source}{}{}{}}
\makeatother
\crefname{originalsection}{original section}{original sections}
\Crefname{originalsection}{Original section}{Original sections}
\makeatletter
\newlabel{algebra-section-orders-of-vanishing}{{02MB}{source}{}{sourcepointer-11}{}}
\newlabel{algebra-section-orders-of-vanishing@cref}{{[originalsection][0][]02MB}{[1][0][]source}{}{}{}}
\makeatother
\title{Chow groups and cycle groups commute with a filtered limit of affine base changes in the c = 0 situation}
\author{The Stacks Project Authors}
\date{stacks-project:chow.tex:local-snapshot-20261012; GNU FDL 1.2 or later}
\begin{document}
\maketitle
\noindent\textit{Editorial note.} All delivered mathematical text is the author's own, copied from the pinned Stacks Project chapter file chow.tex; no statement or proof step was rewritten, reordered or removed.\par\medskip
\noindent\textit{Editorial note.} Preamble provenance: the pinned chapter file opens with an input of the shared preamble and carries no preamble text of its own, so the project's single Stacks-wide preamble staged at sources/stacks-chow/preamble.tex is excerpted verbatim as two spans, lines 14-18 (Xy-pic setup) and lines 154-217 (the native theorem declarations and mathematical macros). Nothing was added, removed or redefined; the wrapper class and the amsmath/amssymb/amsthm/graphicx base packages are supplied by the assembler.\par\medskip
\noindent\textit{Editorial note.} The counted statement and proof contain no cross reference, no citation and no included file, so the delivered body is self contained.\par\medskip
\noindent\textit{Editorial note.} Canonical location of the counted proposition: Stacks tag 0FVQ (sources/stacks-chow/pinned-original-source.tex line 15483).\par\medskip
\noindent\textit{Allowed standard prerequisites (editorial).}\begin{itemize}
\item The cycle groups Z\_k and Chow groups CH\_k attached to the situation, and functoriality under the base changes. Same chapter, Lemma labelled lemma-compose-base-change, cited at proof lines 15500-15503, together with Situation labelled situation-setup-base-change invoked in the statement. Supplies the systems of cycle and Chow groups with their transition maps whose colimits are being computed.
\item Descending finite presentation, closed immersions, opens and sections through a filtered limit of affine schemes. Limits, Lemmas labelled limits-lemma-descend-finite-presentation, limits-lemma-descend-closed-immersion-finite-presentation, limits-lemma-descend-opens and limits-lemma-descend-section, cited at proof lines 15512-15519 and 15562-15566. Produces the schemes, subschemes, opens and rational functions at some index i whose base change is the given object on X'.
\item Behaviour of delta-dimension under base change and nonemptiness in a filtered limit. Same chapter, Lemma labelled lemma-dimension-base-change, cited at proof lines 15524-15526 and 15534-15537; Limits, Lemma labelled limits-lemma-limit-nonempty, cited at proof line 15548. Matches the delta-dimensions of the descended pieces with those on X', and converts emptiness on X' into emptiness at a transition map.
\item Orders of vanishing along a flat local ring extension. Algebra, Lemma labelled algebra-lemma-pullback-module, cited at proof lines 15614-15616, and the construction of ord in Algebra, Section labelled algebra-section-orders-of-vanishing. Gives the equality of the divisor coefficients at the chosen generic points.
\end{itemize}\medskip
\ifdefined\hypertarget\hypertarget{sourcepointer-1}{}\fi
\noindent\textit{Source reference.} Original reference 0FVG: sources/stacks-chow/pinned-original-source.tex, source lines 15123--15123. Same-chapter pointer into the author environment situation at source line 15123; invoked at proof line None of the counted proof. It is retained only as a navigation pointer; every load-bearing step is inside the delivered span.\par\medskip
\ifdefined\hypertarget\hypertarget{sourcepointer-2}{}\fi
\noindent\textit{Source reference.} Original reference 02QL: sources/stacks-chow/pinned-original-source.tex, source lines 1372--1372. Same-chapter pointer into the author environment situation at source line 1372; invoked at proof line None of the counted proof. It is retained only as a navigation pointer; every load-bearing step is inside the delivered span.\par\medskip
\ifdefined\hypertarget\hypertarget{sourcepointer-3}{}\fi
\noindent\textit{Source reference.} Original lemma 0FVP: sources/stacks-chow/pinned-original-source.tex, source lines 15438--15438. Same-chapter pointer into the author environment lemma at source line 15438; invoked at proof line 15501 of the counted proof. It is retained only as a navigation pointer; every load-bearing step is inside the delivered span.\par\medskip
\ifdefined\hypertarget\hypertarget{sourcepointer-4}{}\fi
\noindent\textit{Source reference.} Original lemma 01ZM: sources/stacks-chow/tags.txt, source lines 2599--2599. Cross-chapter standard prerequisite located in the project tag list staged in the local snapshot; invoked at proof line 15513 of the counted proof. It is retained only as a navigation pointer; every load-bearing step is inside the delivered span.\par\medskip
\ifdefined\hypertarget\hypertarget{sourcepointer-5}{}\fi
\noindent\textit{Source reference.} Original lemma 01ZP: sources/stacks-chow/tags.txt, source lines 2602--2602. Cross-chapter standard prerequisite located in the project tag list staged in the local snapshot; invoked at proof line 15517 of the counted proof. It is retained only as a navigation pointer; every load-bearing step is inside the delivered span.\par\medskip
\ifdefined\hypertarget\hypertarget{sourcepointer-6}{}\fi
\noindent\textit{Source reference.} Original lemma 0FVH: sources/stacks-chow/pinned-original-source.tex, source lines 15139--15139. Same-chapter pointer into the author environment lemma at source line 15139; invoked at proof line 15521 of the counted proof. It is retained only as a navigation pointer; every load-bearing step is inside the delivered span.\par\medskip
\ifdefined\hypertarget\hypertarget{sourcepointer-7}{}\fi
\noindent\textit{Source reference.} Original lemma 01Z2: sources/stacks-chow/tags.txt, source lines 2577--2577. Cross-chapter standard prerequisite located in the project tag list staged in the local snapshot; invoked at proof line 15543 of the counted proof. It is retained only as a navigation pointer; every load-bearing step is inside the delivered span.\par\medskip
\ifdefined\hypertarget\hypertarget{sourcepointer-8}{}\fi
\noindent\textit{Source reference.} Original lemma 01Z4: sources/stacks-chow/tags.txt, source lines 2579--2579. Cross-chapter standard prerequisite located in the project tag list staged in the local snapshot; invoked at proof line 15562 of the counted proof. It is retained only as a navigation pointer; every load-bearing step is inside the delivered span.\par\medskip
\ifdefined\hypertarget\hypertarget{sourcepointer-9}{}\fi
\noindent\textit{Source reference.} Original lemma 01Z0: sources/stacks-chow/tags.txt, source lines 2575--2575. Cross-chapter standard prerequisite located in the project tag list staged in the local snapshot; invoked at proof line 15564 of the counted proof. It is retained only as a navigation pointer; every load-bearing step is inside the delivered span.\par\medskip
\ifdefined\hypertarget\hypertarget{sourcepointer-10}{}\fi
\noindent\textit{Source reference.} Original lemma 02M1: sources/stacks-chow/tags.txt, source lines 3428--3428. Cross-chapter standard prerequisite located in the project tag list staged in the local snapshot; invoked at proof line 15615 of the counted proof. It is retained only as a navigation pointer; every load-bearing step is inside the delivered span.\par\medskip
\ifdefined\hypertarget\hypertarget{sourcepointer-11}{}\fi
\noindent\textit{Source reference.} Original section 02MB: sources/stacks-chow/tags.txt, source lines 3438--3438. Cross-chapter standard prerequisite located in the project tag list staged in the local snapshot; invoked at proof line 15617 of the counted proof. It is retained only as a navigation pointer; every load-bearing step is inside the delivered span.\par\medskip
\begin{lemma}
\label{lemma-chow-limit}
In Situation \ref{situation-setup-base-change} assume $c = 0$
and assume that $S' = \lim_{i \in I} S_i$ is a filtered limit
of schemes $S_i$ affine over $S$ such that
\begin{enumerate}
\item with $\delta_i$ equal to $S_i \to S \xrightarrow{\delta} \mathbf{Z}$
the pair $(S_i, \delta_i)$ is as in Situation \ref{situation-setup},
\item $S_i, \delta_i, S, \delta, S \to S_i, c = 0$ is as in
Situation \ref{situation-setup-base-change},
\item $S_i, \delta_i, S_{i'}, \delta_{i'}, S_i \to S_{i'}, c = 0$ 
for $i \geq i'$ is as in Situation \ref{situation-setup-base-change}.
\end{enumerate}
Then for a quasi-compact scheme $X$ of finite type over $S$
with base change $X'$ and $X_i$ by $S' \to S$ and $S_i \to S$ we have
$Z_k(X') = \colim Z_k(X_i)$ and $\CH_k(X') = \colim \CH_k(X_i)$.
\end{lemma}

\begin{proof}
By the result of Lemma \ref{lemma-compose-base-change} we obtain
a system of cycle groups $Z_k(X_i)$ and a system of
chow groups $\CH_k(X_i)$ as well as maps
$\colim Z_k(X_i) \to Z_k(X')$ and $\colim \CH_i(X_i) \to \CH_k(X')$.
We may replace $S$ by a quasi-compact open through which $X \to S$
factors, hence we may and do assume all the schemes occurring in
this proof are Noetherian (and hence quasi-compact and quasi-separated).

\medskip\noindent
Let us show that the map $\colim Z_k(X_i) \to Z_k(X')$ is surjective.
Namely, let $Z' \subset X'$
be an integral closed subscheme of $\delta'$-dimension $k$. By
Limits, Lemma \ref{limits-lemma-descend-finite-presentation}
we can find an $i$ and a morphism $Z_i \to X_i$ of finite presentation
whose base change is $Z'$. After increasing $i$ we may assume $Z_i$
is a closed subscheme of $X_i$, see
Limits, Lemma \ref{limits-lemma-descend-closed-immersion-finite-presentation}.
Then $Z' \to X_i$ factors through $Z_i$ and we may replace $Z_i$
by the scheme theoretic image of $Z' \to X_i$. In this way we see
that we may assume $Z_i$ is an integral closed subscheme of $X_i$.
By Lemma \ref{lemma-dimension-base-change} we conclude that
$\dim_{\delta_i}(Z_i) = \dim_{\delta'}(Z') = k$.
Thus $Z_k(X_i) \to Z_k(X')$ maps $[Z_i]$ to $[Z']$ and
we conclude surjectivity holds.

\medskip\noindent
Let us show that the map $\colim Z_k(X_i) \to Z_k(X')$ is injective.
Let $\alpha_i = \sum n_j[Z_j] \in Z_k(X_i)$ be a cycle whose image
in $Z_k(X')$ is zero. We may and do assume $Z_j \not = Z_{j'}$ if
$j \not = j'$ and $n_j \not = 0$ for all $j$.
Denote $Z'_j \subset X'$ the base change of $Z_j$.
By Lemma \ref{lemma-dimension-base-change} each irreducible component
of $Z'_j$ has $\delta'$-dimension $k$. Moreover, as $Z_j$ is irreducible
and $Z'_j \to Z_j$ is flat (as the base change of $S' \to S$)
we see that $Z'_j \to Z_j$ is dominant.
Hence if $Z'_j$ is nonempty, then some irreducible component, say
$Z'$, of $Z'_j$ dominates $Z_j$. It follows that $Z'$
cannot be an irreducible component of $Z'_{j'}$ for $j' \not = j$.
Hence if $Z'_j$ is nonempty, then we see that $(S' \to S_i)^*\alpha_i
= \sum [Z'_j]_r$ is nonzero (as the coefficient of $Z'$ would be nonzero).
Thus we see that $Z'_j = \emptyset$ for all $j$. However, this means
that the base change of $Z_j$ by some transition map $S_{i'} \to S_i$
is empty by Limits, Lemma \ref{limits-lemma-limit-nonempty}.
Thus $\alpha_i$ dies in the colimit as desired.

\medskip\noindent
The surjectivity of $\colim Z_k(X_i) \to Z_k(X')$ implies that
$\colim \CH_k(X_i) \to \CH_k(X')$ is surjective.
To finish the proof we show that this map is injective.
Let $\alpha_i \in \CH_k(X_i)$
be a cycle whose image $\alpha' \in \CH_k(X')$ is zero.
Then there exist integral closed subschemes
$W_l' \subset X'$, $l = 1, \ldots, r$ of $\delta"$-dimension $k + 1$
and nonzero rational functions $f'_l$ on $W'_l$
such that $\alpha' = \sum_{l = 1, \ldots, r} \text{div}_{W'_l}(f'_l)$.
Arguing as above we can find an $i$ and integral closed subschemes
$W_{i, l} \subset X_i$ of $\delta_i$-dimension $k + 1$
whose base change is $W'_l$.
After increasing $i$ we may assume we have rational functions
$f_{i, l}$ on $W_{i, l}$. Namely, we may think of $f'_l$ as a
section of the structure sheaf over a nonempty open $U'_l \subset W'_l$,
we can descend these opens by Limits, Lemma \ref{limits-lemma-descend-opens}
and after increasing $i$ we may descend $f'_l$ by
Limits, Lemma \ref{limits-lemma-descend-section}.
We claim that
$$
\alpha_i = \sum\nolimits_{l = 1, \ldots, r} \text{div}_{W_{i, l}}(f_{i, l})
$$
after possibly increasing $i$.

\medskip\noindent
To prove the claim, let $Z'_{l, j} \subset W'_l$ be a finite
collection of integral closed subschemes of $\delta'$-dimension $k$
such that $f'_l$ is an invertible regular function outside
$\bigcup_j Y'_{l, j}$. After increasing $i$ (by the arguments above)
we may assume there exist integral closed subschemes $Z_{i, l, j} \subset W_i$
of $\delta_i$-dimension $k$ such that $f_{i, l}$ is an
invertible regular function outside $\bigcup_j Z_{i, l, j}$.
Then we may write
$$
\text{div}_{W'_l}(f'_l) = \sum n_{l, j} [Z'_{l, j}]
$$
and
$$
\text{div}_{W_{i, l}}(f_{i, l}) = \sum n_{i, l, j} [Z_{i, l, j}]
$$
To prove the claim it suffices to show that $n_{l, i} = n_{i, l, j}$.
Namely, this will imply that $\beta_i =
\alpha_i - \sum\nolimits_{l = 1, \ldots, r} \text{div}_{W_{i, l}}(f_{i, l})$
is a cycle on $X_i$ whose pullback to $X'$ is zero as a cycle!
It follows that $\beta_i$ pulls back to zero as a cycle on $X_{i'}$
for some $i' \geq i$ by an easy argument we omit.

\medskip\noindent
To prove the equality $n_{l, i} = n_{i, l, j}$ we choose a
generic point $\xi' \in Z'_{l, j}$ and we denote
$\xi \in Z_{i, l, j}$ the image which is a generic point also.
Then the local ring map
$$
\mathcal{O}_{W_{i, l}, \xi}
\longrightarrow
\mathcal{O}_{W'_l, \xi'}
$$
is flat as $W'_l \to W_{i, l}$ is the base change of the flat
morphism $S' \to S_i$. We also have
$\mathfrak m_\xi \mathcal{O}_{W'_l, \xi'} = \mathfrak m_{\xi'}$
because $Z_{i, l, j}$ pulls back to $Z'_{l, j}$! Thus the equality of
$$
n_{l, j} = \text{ord}_{Z'_{l, j}}(f'_l) =
\text{ord}_{\mathcal{O}_{W'_l, \xi'}}(f'_l)
\quad\text{and}\quad
n_{i, l, j} = \text{ord}_{Z_{i, l, j}}(f_{i, l}) =
\text{ord}_{\mathcal{O}_{W_{i, l}, \xi}}(f_{i, l})
$$
follows from Algebra, Lemma \ref{algebra-lemma-pullback-module}
and the construction of $\text{ord}$ in
Algebra, Section \ref{algebra-section-orders-of-vanishing}.
\end{proof}

\end{document}
